\part{Daily Workflow}

% =====================================================================
\chapter{Getting Started}
\label{ch:start}

\section{What Is LitFolio}

LitFolio is a desktop research paper workbench:
all data---metadata, PDFs, notes, highlights, terms, surveys---lives directly under
\fpath{\textasciitilde/Litera-Library/} on your own machine. Network access is only
needed to fetch arXiv metadata, parse RSS feeds, and call LLMs. In other words, you
can read, write, and search even when offline; moving to another computer is as simple
as syncing that directory.

LitFolio is composed of eight major feature areas:

\begin{enumerate}
  \item \textbf{Library}---local paper list, reading-state machine, TL;DR, deep read,
    translation, folders, tags, automatic BibTeX generation, citation formatting,
    batch export.
  \item \textbf{Reader}---three-pane PDF view with a highlight list on the left and
    notes / selection translation / terms tabs on the right; structured note cards
    and semantic annotation colors.
  \item \textbf{Import}---three entry points: arXiv/DOI metadata, local PDF files,
    and Semantic Scholar search; batch folder import with deduplication.
  \item \textbf{Discovery}---arXiv category browsing, RSS subscriptions, topic survey
    generation, and topic alerts.
  \item \textbf{Ask}---multi-turn RAG conversation over your read papers: query
    rewriting, multi-path FTS5 retrieval, answers with citation numbers, and
    context-aware follow-ups.
  \item \textbf{Knowledge Graph}---paper-relationship network, citation network,
    and concept map, with AI-powered discovery of connections and method evolution.
  \item \textbf{Advanced Tools}---multi-paper comparison, literature review draft,
    similar-paper recommendations, reading queue, smart collections, custom metadata
    fields, Markdown export, command palette.
  \item \textbf{Settings}---multiple LLM endpoint configurations, per-task model
    assignment, whole-library WebDAV sync.
\end{enumerate}

\section{Installation and First Launch}

LitFolio is built with Tauri 2 + React 18 and is available through two channels:

\subsection{Download from Releases (Recommended)}

Go to the project repository's Releases page and download the installer for your platform:

\begin{itemize}
  \item Linux: \fpath{LitFolio\_x.y.z\_amd64.AppImage} or \fpath{.deb}
  \item macOS: \fpath{LitFolio\_x.y.z\_universal.dmg}
  \item Windows: \fpath{LitFolio\_x.y.z\_x64-setup.exe}
\end{itemize}

\begin{tipbox}[Linux AppImage Usage]
After downloading the \fpath{.AppImage} file, make it executable:
\cmd{chmod +x LitFolio\_x.y.z\_amd64.AppImage}, then double-click or run from
the terminal. AppImage requires FUSE (\cmd{libfuse2}); most distributions ship it
by default. If prompted, install it with \cmd{sudo apt install libfuse2}
(Ubuntu/Debian) or \cmd{sudo dnf install fuse-libs} (Fedora).
AppImage is a portable format---no installation needed; simply delete the file to
remove it.
\end{tipbox}

\subsection{Run from Source (Developers)}

\begin{enumerate}
  \item Install \cmd{pnpm}, \cmd{rustup} (with the stable toolchain), and the
    platform-specific Tauri system dependencies
    (on Linux typically \cmd{webkit2gtk-4.1}, \cmd{librsvg}, \cmd{libsoup-3.0}).
  \item Run \cmd{pnpm install} to install frontend dependencies.
  \item Run \cmd{pnpm tauri dev} to start development mode; the first launch will
    compile Rust crates and may take a few minutes.
\end{enumerate}

\begin{tipbox}
In development mode the window title includes a ``[dev]'' badge. The database is
the same as the production build (both live under \fpath{\textasciitilde/Litera-Library/}),
so papers you import, deep-read, or annotate in dev mode will also appear in the
packaged build.
\end{tipbox}

\section{The Library Directory at a Glance}

On first launch, LitFolio automatically creates \fpath{\textasciitilde/Litera-Library/}
and initializes an empty database. You will immediately see the left navigation bar
(\textbf{Library / Import / Browse / Feeds / Topics / Ask / Settings}) and an empty
main list.

\litfigure{01-shell-library-en.png}{LitFolio main interface. Left: navigation bar; center: folder sidebar;
right: paper list with search box. Deep violet (\cmd{violet}) is the primary interactive
color, used for buttons, links, and emphasis.}

The physical layout of the library directory (detailed in Chapter 12) is roughly:

\begin{quote}\small\ttfamily
\textasciitilde/Litera-Library/\\
\ \ library.db\hfill\textnormal{\small // SQLite main database}\\
\ \ papers/<paperId>/\hfill\textnormal{\small // one directory per paper}\\
\ \ \ \ paper.pdf\\
\ \ \ \ meta.json\\
\ \ \ \ note.md\\
\ \ \ \ pdf-text.txt\\
\ \ backups/\hfill\textnormal{\small // automatic backups}\\
\ \ sync/\hfill\textnormal{\small // sync snapshots}
\end{quote}

\section{Next Step: Configure Your First LLM Endpoint}

LitFolio does not ship with any built-in LLM. To use TL;DR, deep reading, translation,
term extraction, topic surveys, and library Q\&A, you must first add at least one
endpoint under \emph{Settings $\to$ LLM Configuration} (any OpenAI-compatible API
or a local Ollama instance will work). See \textbf{Chapter 15} for full details;
you can jump there now and come back, or proceed to Chapter 2 to explore the
features that do not require an LLM.

% =====================================================================
\chapter{The Library}
\label{ch:library}

The library page (\cmd{/library} route) is the screen you will spend the most time on.
Its design goal is simple: one glance at the list should tell you what these papers are
about, which ones deserve a deep read, and which are just RSS items you have not
touched yet.

\section{Main List and Search}

Each row in the main list displays: title, authors, year/journal, reading-status badge,
TL;DR summary (if generated), key-finding chips, and tag chips. Clicking anywhere on
a row opens the detail drawer on the right (see Section 2.4).

\subsection{Search Box}

The search box at the top uses SQLite FTS5: titles, abstracts, TL;DRs, notes, and
highlights are all indexed in a unified full-text index. Keywords are space-delimited
with AND semantics. To search for an exact phrase, wrap it in double quotes.

\subsection{FTS5 Full-Text Search Mechanics}

LitFolio's full-text index is based on the SQLite FTS5 virtual table \cmd{papers\_fts},
using the \cmd{unicode61} tokenizer. The table indexes five fields: title, authors,
abstract, TL;DR, and notes. INSERT / UPDATE / DELETE triggers keep it in sync with
the \cmd{papers} main table---papers are automatically re-indexed on import or update.

The search pipeline:

\begin{enumerate}
  \item \textbf{Input sanitization}---the \cmd{escape\_fts} function splits user input
    on spaces, strips FTS5 special characters (\cmd{"():.-}), wraps each token as
    \cmd{"token"*} (prefix matching), and joins them with AND. For example,
    \cmd{retrieval augmented} becomes \cmd{"retrieval"* AND "augmented"*}.
  \item \textbf{BM25 ranking}---FTS5's built-in BM25 algorithm computes a relevance
    score automatically. Results are sorted by ascending BM25 (lower is more relevant),
    with ties broken by descending import time.
  \item \textbf{Empty-query fallback}---if sanitization produces an empty query (e.g.,
    the user typed only special characters), the system falls back to \cmd{list\_recent},
    returning papers in reverse chronological order.
\end{enumerate}

\begin{tipbox}
The FTS5 \cmd{unicode61} tokenizer splits CJK text character by character, so searching
for a Chinese phrase like ``neural network'' is equivalent to searching for each character
individually with AND---short characters may bring in noise. For precise CJK searching,
consider using English keywords or the English names of technical terms.
\end{tipbox}

\subsection{Sorting and Filtering}

The sidebar's Must-read / Reading / Read / Unread filters correspond to the
\cmd{must}/\cmd{reading}/\cmd{read}/\cmd{unread} status values. Clicking a folder
(see 2.6) or a tag chip applies the corresponding filter.

\section{Reading States: unread → reading → read → must}

Each paper has a \textbf{reading state}, initially \cmd{unread}:

\begin{description}
  \item[unread] Not yet touched. Papers imported via RSS or arXiv default to this state.
  \item[reading] Currently being read. Signals that you want this paper to appear
    near the top of the list.
  \item[read] Finished; the paper can fade into the background.
  \item[must] ``Must read''---a higher-priority marker than \cmd{read}, also given
    priority when generating topic surveys.
\end{description}

Clicking the status badge cycles through the states in order. This four-state mechanism
is a deliberate design choice: LitFolio does not assume that reading a paper once means
you are done with it---it preserves the semantics of ``important papers deserve to be
revisited.''

\section{Quick Read (TL;DR)}

The main list has a ``Quick Read'' button on the right. Clicking it invokes the
``TL;DR generation'' task (bound by default to your configured \cmd{tldr} model).
The returned paragraph is saved directly to the database as a compact summary and
appears instantly the next time you open the paper.

\begin{tipbox}
Once generated, the TL;DR is cached in SQLite. Clicking the same paper again will not
re-invoke the LLM. To force regeneration, click ``Regenerate'' in the drawer.
\end{tipbox}

\section{Deep Read: Four-Section Drawer}

``Deep Read'' is one of the core interactions that sets LitFolio apart from other
reader tools. Click the ``Deep Read'' button on a paper row and the right drawer
presents content in four structured sections:

\litfigure{02-library-drawer-en.png}{Deep-read drawer. Top to bottom: \textbf{What problem
does it solve}, \textbf{What method is proposed}, \textbf{How does it differ from
others}, \textbf{Limitations and open issues}---each section is extracted from the
full text and has a \cmd{Regenerate} button beside it.}

The four-section design lets you judge whether a paper deserves a detailed read
without going through the entire text. The \textbf{Differences} and \textbf{Limitations}
sections alone are usually enough to decide whether to open the PDF. The full result
is likewise cached.

\section{Title/Abstract Translation}

The drawer includes a ``Translate'' button that calls the ``translation'' task (bound
to your \cmd{translate} model). The result is presented as a bilingual side-by-side
view, and the translated title and abstract are saved in the metadata for quick
scanning.

\begin{tipbox}
Different tasks can be bound to different models---for instance, you can assign TL;DR
to DeepSeek and translation to a local Ollama instance. See \textbf{Chapter 15}.
\end{tipbox}

\section{Folders}

Folders are hierarchical (unlike tags) and support multiple membership---a single paper
can reside in both ``ML/Transformer'' and ``Applications/NLP'' simultaneously.

\litfigure{03-library-folder-en.png}{Sidebar folders. Selecting a folder filters the main
list to its contents; combined with the search box, this gives you ``full-text search
within a folder.''}

To create a folder: right-click an empty area in the sidebar $\to$ \cmd{New}. To manage
membership: drag and drop, or click ``Move to\ldots'' in the drawer.

\section{Tags}

Tags are flat, string-based labels created through the input box in the drawer. A paper
can have any number of tags. The distinction from folders:

\begin{itemize}
  \item Folders = \textbf{research domain taxonomy}, relatively stable, 2--3 levels deep.
  \item Tags = \textbf{ad-hoc topic/project markers}, ephemeral and flat.
\end{itemize}

\section{Deleting a Paper}

The drawer has a ``Delete'' button at the bottom. Deleting removes: the entire
\fpath{papers/<id>/} directory, the SQLite row, and all associated highlights,
notes, and terms.

\begin{warnbox}
Deletion is irreversible in the database, but LitFolio moves \fpath{papers/<id>/} to
\fpath{backups/deleted/<timestamp>/} instead of running \cmd{rm -rf}. If you
accidentally delete something, you can recover it from \fpath{backups/deleted/}.
\end{warnbox}

\section{Automatic BibTeX Generation}

When a paper is imported, LitFolio automatically generates a BibTeX entry from its
metadata and stores it in the \cmd{bibtex} database field. Both the drawer and the
reader header have a ``Copy BibTeX'' button that copies the entry to the system
clipboard.

\litfigure{21-library-bibtex-en.png}{BibTeX section in the paper drawer. One click copies the
entry to the clipboard, ready to paste into a LaTeX document.}

\subsection{BibTeX Key Generation Algorithm}

The BibTeX key is a paper's unique citation identifier, formatted as
\cmd{<first-author-surname-lowercase><year><first-significant-title-word-lowercase>},
e.g.\ \cmd{vaswani2017attention}. The generation pipeline:

\begin{enumerate}
  \item \textbf{Extract first author's surname}---take the first element of the
    \cmd{authors} array, split on space or comma, take the last token (to handle both
    \cmd{Firstname Lastname} and \cmd{Lastname, Firstname} formats), lowercase it,
    and strip non-ASCII characters (keeping only \cmd{a-z} and hyphens). If the
    author list is empty, use \cmd{unknown}.
  \item \textbf{Extract first significant title word}---split the title on spaces,
    skip stop words (\cmd{a/an/the/on/in/of/for/and/or/to/with/from/by/at}),
    take the first non-stop word, lowercase it, and keep only \cmd{a-z0-9}. If the
    title consists entirely of stop words (extremely rare), take the first word.
  \item \textbf{Concatenate}---\cmd{<surname><year><first-word>}. If the year is
    missing, use \cmd{0000}.
  \item \textbf{Deduplicate}---after generation, check whether the key already exists
    in the library. If so, append \cmd{a}, \cmd{b}\ldots until the key is unique.
    This ensures that importing the same paper from different sources will never
    produce a BibTeX key collision.
\end{enumerate}

\subsection{Field Mapping}

\begin{longtable}{p{3.5cm} p{3cm} p{6cm}}
\toprule
\textbf{BibTeX Field} & \textbf{Source} & \textbf{Processing Rule} \\
\midrule
\endhead
\cmd{title} & \cmd{papers.title} & Used directly \\
\cmd{author} & \cmd{papers.authors} & JSON array joined with \cmd{ and } \\
\cmd{year} & \cmd{papers.year} & Used directly; omitted if missing \\
\cmd{journal}/\cmd{booktitle} & \cmd{papers.venue} & Written if venue exists; omitted otherwise \\
\cmd{doi} & \cmd{papers.doi} & Written if DOI exists; omitted otherwise \\
\cmd{eprint} & \cmd{papers.arxiv\_id} & Written if arXiv ID exists \\
\cmd{archivePrefix} & Hardcoded & Fixed to \cmd{arXiv} when arXiv ID exists \\
\cmd{primaryClass} & arXiv ID parsing & Inferred from the arXiv ID prefix (e.g.\
  \cmd{cs} $\to$ \cmd{cs.AI}); omitted if uninferrable \\
\cmd{url} & Composed & If DOI exists $\to$ \cmd{https://doi.org/<doi>}; if arXiv ID
  exists $\to$ \cmd{https://arxiv.org/abs/<id>}; otherwise omitted \\
\bottomrule
\end{longtable}

\subsection{Generation on Import}

BibTeX generation is the final step of every import path: arXiv ID import, DOI import,
PDF file import, and Semantic Scholar search import. It is a pure in-memory string
concatenation---no LLM calls, no network access, no external service dependency;
the time cost is negligible. Legacy papers with a NULL BibTeX field can be backfilled
in bulk via the ``Backfill BibTeX'' button on the Settings page.

\begin{tipbox}
BibTeX key uniqueness is enforced by the database. If you use Zotero or another tool
that generates keys in a different format, LitFolio's keys will not collide---because
LitFolio keys always start with a lowercase surname and never contain underscores.
Generation is pure in-memory string concatenation with negligible overhead.
\end{tipbox}

\section{Citation Formatting}

In addition to BibTeX, LitFolio can produce formatted citations in standard academic
styles. In the paper drawer, click ``Copy Citation'' and choose from four styles.
Formatting is pure template rendering---metadata fields are substituted into predefined
string templates; no LLM is involved.

\subsection{Four Style Templates}

\subsubsection{APA 7th}

Template:
\cmd{<Surname>, <First Initial>. (<Year>). <Title>. <Journal>.}

Example output:\\
Vaswani, A., Shazeer, N., Parmar, N., Uszkoreit, J., Jones, L., Gomez, A. N.,
Kaiser, L., \& Polosukhin, I. (2017). Attention is all you need. \textit{NeurIPS}.

Rules: for more than 20 authors, list the first 19 followed by \cmd{...} and the last;
title in sentence case; journal name in italics.

\subsubsection{IEEE}

Template:
\cmd{[N] <First Initial>. <Surname> et al., ``<Title>,'' <Journal>, <Year>.}

Example output:\\
{[}1{]} A. Vaswani et al., ``Attention is all you need,'' \textit{NeurIPS}, 2017.

Rules: IEEE uses bracketed numbering; 3+ authors use \cmd{et al.}; title in Title Case;
journal in italics.

\subsubsection{GB/T 7714-2015}

Template:
\cmd{[N] <SURNAME> <First Initial>, <SURNAME> <First Initial>, ... <Title>[J]. <Journal>, <Year>.}

Example output:\\
{[}1{]} VASWANI A, SHAZEER N, PARMAR N, et al. Attention is all you need[J].
NeurIPS, 2017.

Rules: Chinese national standard---surname in uppercase before initials; 3+ authors
use \cmd{et al.}; document type identifier \cmd{[J]} for journals, \cmd{[C]} for
conferences, \cmd{[D]} for dissertations. LitFolio selects \cmd{[J]} or \cmd{[C]}
automatically based on whether the venue string contains
\cmd{Conference/Proceedings/Workshop}.

\subsubsection{Chicago 17th}

Template:
\cmd{<Surname>, <First>. <Year>. ``<Title>.'' <Journal>.}

Example output:\\
Vaswani, Ashish, Noam Shazeer, Niki Parmar, et al. 2017.
``Attention Is All You Need.'' \textit{NeurIPS}.

Rules: authors use full names (not abbreviations); 7+ authors list the first 7
followed by \cmd{et al.}

\subsection{Batch Citation Formatting}

Select multiple papers and click ``Export Citations $\to$ Formatted Text''. LitFolio
generates a numbered reference list in the chosen style. Numbers are assigned in the
order papers appear in the list (IEEE and GB/T 7714 require numbering; APA and Chicago
sort by author-year).

\section{Batch Citation Export}

In the library, select papers in checkbox mode; a toolbar ``Export Citations'' button
appears. Three formats are supported:

\subsection{BibTeX (.bib)}

Merges the BibTeX entries of all selected papers into a single \cmd{.bib} file. Each
record is taken directly from \cmd{papers.bibtex} (generated by the algorithm in
Section 2.9) and separated by blank lines. If a paper's BibTeX field is NULL (very
old data), the entry is skipped and listed in the export report. The output file can
be used directly with \cmd{bibtex} / \cmd{biber} / \cmd{biblatex}.

\subsection{RIS (.ris)}

RIS (Research Information Systems) is a universal reference exchange format. Each
record consists of tagged lines:

\begin{verbatim}
TY  - JOUR
TI  - Attention Is All You Need
AU  - Vaswani, Ashish
AU  - Shazeer, Noam
PY  - 2017
JO  - NeurIPS
DO  - 10.48550/arXiv.1706.03762
ER  -
\end{verbatim}

Tag mapping rules: \cmd{TY} is determined by venue (journal $\to$ \cmd{JOUR},
conference $\to$ \cmd{CONF}, preprint $\to$ \cmd{UNPB}); \cmd{AU} one author per
line in \cmd{Lastname, Firstname} format; \cmd{DO} from DOI; \cmd{UR} from URL.
RIS files can be imported directly into Zotero, Mendeley, or EndNote.

\subsection{Formatted Text}

Generates a numbered citation list in one of the four styles (APA/IEEE/GB/T 7714/Chicago)
as described in Section 2.10. Each paper occupies one or more lines (depending on
style requirements); numbers are assigned in list-appearance order.

\subsection{File Naming and Trigger}

Export files are named \cmd{litera-export-<YYYY-MM-DD>.<ext>} and saved via a browser
download dialog to a location of your choice. You can also choose ``Copy to Clipboard''
to paste directly into an editor.

\section{Smart Collections}

Smart collections are rule-based virtual folders---papers are automatically included
based on rules, without manual drag-and-drop. The sidebar's ``Smart Collections''
area lets you create and manage them.

\litfigure{22-smart-collections-en.png}{Smart collection creation dialog. Papers are filtered
by combining conditions; results update in real time.}

\subsection{Rule Structure}

Each smart collection stores its rules as a JSON tree supporting nested AND/OR logic:

\begin{verbatim}
{
  "op": "and",
  "conditions": [
    { "field": "read_status", "op": "eq", "value": "unread" },
    { "field": "year", "op": "gte", "value": 2024 },
    { "field": "tags", "op": "includes", "value": "transformer" },
    { "op": "or", "conditions": [
      { "field": "folder", "op": "in", "value": "ML" },
      { "field": "title", "op": "contains", "value": "attention" }
    ]}
  ]
}
\end{verbatim}

\subsection{SQL Translation Engine}

At query time, the rule tree is recursively translated into a SQL WHERE clause:

\begin{enumerate}
  \item \textbf{Leaf nodes}---each condition becomes a SQL fragment:
    \begin{itemize}
      \item \cmd{read\_status} $\to$ \cmd{p.read\_status = '<value>'}
      \item \cmd{year} + \cmd{gte/lte} $\to$ \cmd{p.year >= <value>} or
        \cmd{p.year <= <value>}
      \item \cmd{tags} + \cmd{includes} $\to$ \cmd{EXISTS (SELECT 1 FROM paper\_tags pt
        JOIN tags t ON pt.tag\_id = t.id WHERE pt.paper\_id = p.id AND t.name = '<value>')}
      \item \cmd{folder} + \cmd{in} $\to$ \cmd{EXISTS (SELECT 1 FROM paper\_folders pf
        JOIN folders f ON pf.folder\_id = f.id WHERE pf.paper\_id = p.id
        AND f.name LIKE '<value>\%')} (prefix matching to include subfolders)
      \item \cmd{title} + \cmd{contains} $\to$ \cmd{p.title LIKE '\%<value>\%'}
      \item \cmd{authors} + \cmd{contains} $\to$ \cmd{p.authors LIKE '\%<value>\%'}
    \end{itemize}
  \item \textbf{Branch nodes}---\cmd{op: "and"} becomes \cmd{(cond1 AND cond2 AND \ldots)},
    \cmd{op: "or"} becomes \cmd{(cond1 OR cond2 OR \ldots)}. Nesting depth is unlimited.
  \item \textbf{Final query}---assembled as \cmd{SELECT p.* FROM papers p WHERE <rules>
    ORDER BY p.added\_at DESC}.
\end{enumerate}

\subsection{Real-Time Updates}

Smart collections are not snapshots---they query live every time you select them. So
when a paper's tags or status change, the next time you click the collection the results
immediately reflect the update. This is fundamentally different from ordinary folders:
folders require manual drag-in/drag-out; smart collections only care whether the rules
match.

\subsection{Performance}

The translated SQL leverages SQLite indexes. For tag and folder conditions, the
\cmd{paper\_tags} and \cmd{paper\_folders} tables have \cmd{paper\_id} indexes, so
EXISTS subqueries are fast. Ten smart collections add no perceptible delay to sidebar
loading---only the currently selected one executes its query; the rest display only
their name and rule description.

\begin{tipbox}
Smart collections are ideal for long-lived perspectives such as ``unread ML papers
from the last two years'' or ``all must-read papers.'' For ad-hoc filtering, the search
box is more convenient. Rules can be edited at any time; matching results update
immediately after modification.
\end{tipbox}

\section{Reading Queue}

The reading queue is a prioritization tool independent of folders and tags, solving
the problem of ``I have 20 papers I want to read and need to order them.'' Access it
from the library sidebar under ``Reading Queue.''

\litfigure{23-reading-queue-en.png}{Reading queue. Drag to reorder, set target dates,
overdue papers are highlighted.}

\subsection{Queue Data Model}

The queue is stored in the \cmd{reading\_queue} table. Each row contains:

\begin{itemize}
  \item \cmd{paper\_id}---the associated paper (primary key; a paper can appear in the
    queue only once).
  \item \cmd{priority}---integer priority determined by drag-sort position (starting
    at 0, incrementing).
  \item \cmd{target\_date}---optional Unix timestamp representing the planned reading
    date.
  \item \cmd{note}---optional note text.
  \item \cmd{added\_at}---timestamp when the paper was enqueued.
\end{itemize}

\subsection{Drag-and-Drop Sorting}

The frontend uses the HTML5 Drag and Drop API for sorting. When a drag completes:

\begin{enumerate}
  \item The new position index of the dragged item is computed.
  \item All queue items have their \cmd{priority} values reassigned (starting at 0,
    incrementing by 1).
  \item The \cmd{queue\_reorder} command batch-updates the database.
\end{enumerate}

Sort order persists across sessions---close the app and reopen it, and the order is
unchanged.

\subsection{Target Dates and Overdue Alerts}

Target dates are optional. If the current date exceeds a paper's target date, that row
is highlighted in amber (the same color used by \cmd{warnbox}). Overdue status is a
visual cue only; it does not change any paper state automatically.

\subsection{Automatic Status Update}

When you open a queued paper in the reader, the system automatically transitions that
paper's \cmd{read\_status} from \cmd{unread} to \cmd{reading}. This transition is
idempotent---if the paper is already \cmd{reading} or higher, nothing happens.

\subsection{Queue vs.\ Folders/Tags}

The queue is a completely independent dimension:

\begin{itemize}
  \item A paper can be in any folder, carry any tag, and be in any reading state while
    also being in the queue.
  \item Removing a paper from the queue does not delete the paper itself or change its
    folders, tags, or status.
  \item If a queued paper is deleted, its queue entry is cascade-deleted.
\end{itemize}

\section{Paper Deduplication}

LitFolio automatically detects duplicate papers at import time, preventing the same
paper from appearing multiple times due to different import sources (arXiv ID / DOI /
PDF file).

\subsection{Three-Level Matching Strategy}

Detection runs in priority order and stops on first match:

\begin{enumerate}
  \item \textbf{DOI exact match}---query \cmd{SELECT id FROM papers WHERE doi = ?}.
    DOIs are globally unique identifiers with a 100\% match rate.
  \item \textbf{arXiv ID exact match}---query \cmd{SELECT id FROM papers WHERE
    arxiv\_id = ?}. Different versions of the same paper (\cmd{v1}/\cmd{v2}) have
    their version suffix stripped before comparison.
  \item \textbf{Title similarity match}---uses normalized Levenshtein edit distance.
    Given two titles $s_1, s_2$, the edit distance $d(s_1, s_2)$ (each insertion,
    deletion, or substitution costs 1) is normalized as:
    \[
      D(s_1, s_2) = \frac{d(s_1, s_2)}{\max(|s_1|, |s_2|)}
    \]
    When $D < 0.15$, the papers are considered duplicates. This means two titles can
    differ by up to 15\% of their characters---enough to tolerate case differences
    (\cmd{Attention Is All You Need} vs.\ \cmd{attention is all you need}),
    punctuation differences (hyphens vs.\ spaces), and minor OCR errors.
\end{enumerate}

\subsection{Title Preprocessing}

Before computing edit distance, both titles undergo preprocessing:

\begin{enumerate}
  \item Convert to lowercase.
  \item Remove all non-alphanumeric characters (keeping spaces).
  \item Collapse consecutive spaces into a single space.
  \item Trim leading and trailing whitespace.
\end{enumerate}

This ensures that \cmd{"Attention-Is-All-You-Need"} and
\cmd{"attention is all you need"} have a distance of 0 (exact match).

\subsection{Full-Library Scan}

The ``Find Duplicates'' button on the Settings page performs a full-library scan.
The algorithm is $O(n^2)$ title comparisons (where $n$ is the number of papers),
with two optimizations:

\begin{itemize}
  \item \textbf{DOI/arXiv ID pre-filter}---papers with DOIs are bucketed by DOI first;
    title comparison only happens within the same bucket. The same applies to arXiv
    IDs. This dramatically reduces the number of pairs needing string comparison.
  \item \textbf{Length pre-filter}---if the length difference between two titles exceeds
    15\%, the edit-distance computation is skipped entirely. This is a conservative
    pruning: when $|s_1| - |s_2| > 0.15 \times \max(|s_1|, |s_2|)$, $D$ is
    necessarily $> 0.15$.
\end{itemize}

A full scan of 1,000 papers takes about 2--3 seconds on typical hardware.

\subsection{Merge Workflow}

When duplicates are found, a merge dialog appears. You choose which record to keep
(usually the one with more complete metadata), and the system performs:

\begin{enumerate}
  \item Transfer all highlights, notes, terms, tags, and folder memberships from the
    merged paper to the retained paper.
  \item Set the retained paper's reading state to the higher of the two
    (\cmd{must} > \cmd{read} > \cmd{reading} > \cmd{unread}).
  \item Fill empty TL;DR, deep-read, and translation fields on the retained paper
    with values from the merged paper (existing values are not overwritten).
  \item Delete the merged paper's database row and \fpath{papers/<id>/} directory.
\end{enumerate}

\begin{warnbox}
Merge operations cannot be undone. The merged paper's PDF file is moved to
\fpath{backups/deleted/}, but metadata is deleted directly. If you are unsure which
record to keep, cancel the merge, manually inspect both records, and try again.
\end{warnbox}

% =====================================================================
\chapter{Importing Papers}
\label{ch:import}

The Import page (\cmd{/import}) consolidates three entry points: arXiv/DOI metadata,
local PDFs, and Semantic Scholar search. All three ultimately write to a
\fpath{papers/<id>/} directory and a SQLite row.

\section{arXiv ID / DOI}

The most common import path. It is a \textbf{two-step} process:

\litfigure{04-import-arxiv-en.png}{arXiv/DOI tab. Paste an ID or DOI, click ``Query Metadata''
to retrieve title/authors/abstract, then click ``Import'' to persist.}

\begin{enumerate}
  \item \textbf{Query metadata}---paste an arXiv ID like \cmd{1706.03762} or a DOI like
    \cmd{10.1038/s41586-024-...} and click Query. LitFolio fetches metadata from the
    arXiv API or Crossref and immediately displays the title, authors, and abstract.
  \item \textbf{Import}---confirm the metadata and click ``Import.'' LitFolio will
    simultaneously attempt to download the PDF (directly from arXiv for arXiv IDs;
    via Unpaywall's open-access pathway for DOIs---a download failure does not block
    import).
\end{enumerate}

\begin{tipbox}
arXiv IDs support both the new format (\cmd{2402.12345}) and the old format
(\cmd{cs/0501001}). Versioned IDs like \cmd{2402.12345v2} are also parsed correctly.
\end{tipbox}

\section{Local PDF Files}

\litfigure{05-import-pdf-en.png}{PDF file tab. Select a single PDF or an entire directory;
edit the guessed title/author before importing.}

\begin{itemize}
  \item \textbf{Single PDF}---after selecting a file, LitFolio runs lightweight metadata
    extraction (first page of the PDF + DOI regex) to guess title and authors; you can
    edit them in the form.
  \item \textbf{Drag and drop}---drag a PDF onto the LitFolio window with
    \textbf{targeted drop zones}: dropping onto a paper row in the library list
    directly binds the PDF (replacing manual file selection); dropping onto the
    PDF area on the import page fills the selected file; dropping onto the DOI
    citation import dialog fills the PDF. If no specific target is hit, the
    default batch import flow is used.
\end{itemize}

\subsection{PDF Metadata Extraction Strategy}

When importing a local PDF, LitFolio tries the following extraction steps in order:

\begin{enumerate}
  \item \textbf{PDF Info Dictionary}---use \cmd{lopdf} to read the \cmd{/Title},
    \cmd{/Author}, and \cmd{/Subject} fields. Text decoding attempts both UTF-8 and
    UTF-16 BE (with BOM), since academic PDFs commonly use either encoding.
  \item \textbf{First-page DOI regex}---if the Info Dictionary contains no DOI, scan
    the first-page text with the regex \cmd{10.\textbackslash d\{4,9\}/\textbackslash S+},
    stripping trailing punctuation (.,);] etc.) from matches.
  \item \textbf{Title guess}---if still no title, take the first line on page 1 that
    is $\geq 12$ characters long and contains letters.
  \item \textbf{Filename fallback}---if everything else fails, use the PDF filename
    (without extension) as the title.
\end{enumerate}

The author field is split on commas, semicolons, and ``and'' (accommodating both
BibTeX and natural-language formats). A SHA-256 hash of the file is also computed
at import time for deduplication.

\section{Semantic Scholar Search}

If you only remember a fragment of the title, a keyword, or want to find all papers
by a specific author:

\litfigure{06-import-search-en.png}{Search tab. Results are sorted by citation count;
each can be directly ``Imported'' (fetches metadata + attempts PDF download).}

LitFolio calls the Semantic Scholar Graph API, returning 25 results per page. After
selecting a result and clicking ``Import,'' the full metadata is fetched and the PDF
download is attempted.

\begin{warnbox}
The Semantic Scholar public API has QPS limits; frequent searching may trigger 429
responses. If your research direction requires heavy retrieval, consider applying for
a free API key on their website and entering it under ``Settings $\to$ API Keys.''
\end{warnbox}

\section{Batch Folder Import}

Beyond single-PDF import, LitFolio supports importing all PDFs in a directory. On the
Import page, click the ``Import Folder'' button and select a directory.

\subsection{Recursive Scanning}

After selecting a folder, LitFolio recursively scans all subdirectories for \cmd{.pdf}
files. The scan results are displayed as a list, each entry showing:

\begin{itemize}
  \item Filename (used as a fallback source for title guessing).
  \item File size.
  \item Extraction status (Pending / Processing / Success / Failed).
\end{itemize}

\subsection{Filename Heuristic Parsing}

For batch-imported PDFs, LitFolio attempts to extract metadata from the filename.
Common naming patterns:

\begin{description}
  \item[\cmd{Author\_Year\_Title.pdf}] Split on underscores: first segment is the
    author, second is the year (4 digits), the rest is the title.
  \item[\cmd{Author - Title (Year).pdf}] Split on \cmd{ - } to separate author and
    title; extract the year from parentheses.
  \item[\cmd{2024 - Author - Title.pdf}] Year-first variant.
\end{description}

Heuristic parsing achieves roughly 70\% accuracy. For files that cannot be parsed,
the filename (minus extension) is used as the title, and the author and year are left
blank for manual completion after import.

\subsection{Concurrent Processing and Progress Tracking}

Batch import uses 4-way concurrent processing (configurable). Progress is pushed to
the frontend in real time via event streams:

\begin{itemize}
  \item A progress bar shows ``Processing 23/150\ldots\ (3 failed).''
  \item Failed files are listed separately with failure reasons (corrupted PDF, no text
    layer, metadata extraction failure, etc.).
  \item On completion, a summary is displayed: N succeeded, N failed, N skipped
    (duplicate papers already in the library).
\end{itemize}

\begin{tipbox}
Batch-importing 100 PDFs typically completes within 5 minutes. If some PDFs are
scanned images (no text layer), LitFolio will still import them (using the filename
as the title), but full-text search and RAG retrieval will not cover their content.
You can preprocess scanned PDFs with \cmd{ocrmypdf} and then re-bind them.
\end{tipbox}

\section{Redirect from RSS / arXiv Browse}

The arXiv browse and RSS feed features described in Chapters 5 and 6 each include an
``Import'' button in their detail drawers, which is equivalent to pasting the arXiv ID
on the Import page. In other words: you do not need to visit the Import page
separately---when browsing RSS, just import directly from what catches your eye.

% =====================================================================
\chapter{The Reader}
\label{ch:reader}

The reader is LitFolio's second most-used page. Access it from the main list by
clicking the \kbd{clip} icon or ``Open in Reader.''

\section{Three-Pane Layout}
\subsection{PDF Text Extraction and Caching}

When the reader loads a PDF, PDF.js performs full-text extraction on the frontend
(it handles CMap encoding, embedded font subsets, and ToUnicode tables in modern
academic PDFs more reliably than the backend \cmd{lopdf}). The extracted text is
sent to the backend via IPC and cached in \fpath{papers/<id>/text.txt}. Subsequent
term extraction and FTS5 indexing read from this cache file directly---as long as
\fpath{text.txt}'s modification time is newer than \fpath{paper.pdf}'s, the cache
is reused; otherwise \cmd{lopdf} re-extracts. This means term generation never
needs to re-parse the PDF and is available instantly.

\litfigure{14-reader-en.png}{Reader three-pane layout. Left: highlight list; center:
PDF.js-rendered document; right: notes / selection translation / terms tab switcher.
Pane widths are adjustable by dragging; state is automatically persisted.}

\begin{itemize}
  \item \textbf{Left pane}---all highlights for the current paper (sorted by page);
    click any highlight to jump to its location.
  \item \textbf{Center pane}---the PDF body. Based on PDF.js + react-pdf-highlighter;
    inherits system dark mode and inverts colors during rendering. A zoom toolbar in
    the top-left corner supports the following operations:
    \begin{itemize}
      \item \kbd{Ctrl++} / \kbd{Ctrl+-}: zoom in / out in 15\% increments.
      \item \kbd{Ctrl+0}: reset to 100\%.
      \item \kbd{Ctrl+scroll wheel}: smooth zoom.
      \item The ``Fit'' button (\cmd{RotateCcw} icon) restores page-width auto-fit.
    \end{itemize}
    Zoom range is 60\%--240\%, defaulting to ``page-width'' fit mode.
  \item \textbf{Right pane}---three online tool tabs: \textbf{Notes} (Markdown),
    \textbf{Selection Translation}, and \textbf{Terms}.
\end{itemize}

\section{Highlights: Text Selection + Region}

Two highlight methods:

\begin{description}
  \item[Text selection highlight] Select text in the PDF as you normally would;
    releasing the mouse shows a small popup---click ``Highlight.'' The selected text
    is cleaned of ligatures and line breaks and written to the highlights table.
  \item[Region highlight (\kbd{Alt+drag})] Hold Alt and drag a rectangle over the PDF
    to capture a screenshot as a highlight---useful for formulas, charts, and diagrams.
\end{description}

Each highlight in the left pane supports comments (click the pencil icon); comments
are written in Markdown.

\section{DOI Citation Import}

When reading a PDF, if the text contains clickable DOI links (e.g.,
\cmd{https://doi.org/10.xxxx/xxxx}), LitFolio intercepts the click and opens a
\textbf{citation import dialog} instead of navigating to an external browser.
The dialog automatically checks whether the DOI already exists in the library:

\begin{itemize}
  \item \textbf{Already exists}---shows the paper's title, authors, and year. Click
    ``Create link'' to establish a \cmd{builds_on} citation relationship between the
    current paper and the target (automatically deduplicated).
  \item \textbf{Not found}---shows metadata preview parsed from the DOI. You need to
    select or drop the corresponding PDF file, then click ``Import and link'' to
    simultaneously import the paper and create the citation link.
  \item \textbf{Points to current paper}---shows a notice and disables the action button.
\end{itemize}

\begin{tipbox}
Citation relationships are visualizable in the knowledge graph. You can also drag and
drop a PDF onto a paper row in the library list to quickly bind a PDF file.
\end{tipbox}

\section{Selection Translation}

Select text in the PDF, switch to the ``Selection Translation'' tab in the right pane,
and click ``Translate Selection'':

\litfigure{15-reader-translate-en.png}{Selection translation. In addition to the bilingual
side-by-side view, domain terms from the passage are listed as purple chips---click
any chip to jump to the ``Terms'' tab for the full definition.}

The output is structured as three parts: bilingual text, term extraction, and
cross-paper references.

\subsubsection{Internal Pipeline}

Selection translation is not a single API call but a three-step process:

\begin{enumerate}
  \item \textbf{Term matching}---the selected text is checked against the current
    paper's existing terms (case-insensitive substring matching on normalized forms).
    Matches reuse existing definitions and evidence snippets; misses trigger real-time
    extraction: three regex patterns capture all-uppercase abbreviations (e.g.\ SSIM),
    Title Case proper nouns (e.g.\ MetaDesigner), and 1--3 word noun phrases, filtered
    through \cmd{is\_term\_candidate}.
  \item \textbf{Cross-paper association}---for each term, an FTS5 search across the
    entire library (max 12 results) is performed. After excluding the current paper,
    up to 3 related papers are retained. Association type is determined by which section
    of the target paper the term appears in: method section $\to$ ``Method
    Implementation,'' comparison section $\to$ ``Comparative Discussion,'' other
    $\to$ ``Related Mention.''
  \item \textbf{LLM translation}---the selected text (truncated to 2000 characters) is
    sent to the translation model. The prompt injects a term glossary (matched terms
    with their local definitions) and instructs the LLM to keep English terms verbatim,
    annotating them on first appearance as \cmd{Term[Chinese gloss]}. This keeps the
    translation fluent while preserving term anchors.
\end{enumerate}

The frontend composites the three results: bilingual text + purple term-chip list +
cross-paper links.
\begin{itemize}
  \item \textbf{Bilingual text}---original and translated paragraphs rendered
    side-by-side.
  \item \textbf{Terms}---domain terms extracted from the passage, each with a one-line
    in-context definition; if other papers also use the same term, a ``cross-paper
    reuse \cmd{N}'' badge is shown.
  \item \textbf{Cross-paper references}---locations in the current paper where the same
    concept was previously discussed (page + line); useful for the ``Have I seen this
    before?'' scenario in long papers.
\end{itemize}

\section{Term Table and Term Network}

\litfigure{16-reader-terms-en.png}{Terms tab. After generating the paper's term table, each
term shows ``In-paper evidence / Cross-paper reuse'' sections, facilitating cross-paper
consistency checks.}

``Generate Terms'' executes in \textbf{two phases}:

\begin{enumerate}
  \item \textbf{Candidate extraction}---local TF-IDF statistical filtering with no LLM
    calls, completing instantly. The candidate term list appears immediately, with each
    term's definition area showing ``Waiting for model explanation'' and a spinner icon.
  \item \textbf{Definition generation}---candidates are sent to the LLM for one-line
    definitions. A progress indicator at the top displays ``Candidate terms are ready.
    Generating explanations\ldots''
\end{enumerate}

This approach lets you review and edit the candidate list (removing unwanted terms)
without waiting for the LLM to finish. Generated terms are cached in the
\fpath{papers/<id>/} directory and the SQLite \cmd{terms} table for reuse during
selection translation.

\subsection{Term Extraction Algorithm}

Term generation does not simply dump the full text into an LLM and ask it to list
terms. LitFolio uses a \textbf{local statistical filtering + LLM definition} two-stage
architecture: first, an offline TF-IDF statistical method extracts candidate terms
from the full text; then the LLM writes a one-line definition for each candidate.
The advantages: it saves tokens (no need for the LLM to scan the full text) and the
candidate list is reusable for term matching during selection translation.

\subsubsection{Candidate Extraction: Weighted TF x IDF x Multi-Factor Scoring}

The candidate pipeline has four steps:

\begin{enumerate}
  \item \textbf{Weighted segmentation}---different sections of the paper are assigned
    different weights (title 3.0, abstract 2.0, methods 2.0, research questions 1.5,
    datasets 1.5, comparison 1.0, limitations 1.0, TL;DR 1.0, full PDF text 0.4).
    Terms appearing in the title or abstract carry higher weight; those in the body
    carry lower weight.
  \item \textbf{Regex extraction + filtering}---regex matching is applied to each
    segment: by default only single tokens are captured (proper nouns, all-uppercase
    abbreviations, hyphenated compounds); for ``Full Name (ACRO)'' abbreviation pairs,
    a dedicated recognition path allows the full name up to 5 words. Each candidate
    must pass \cmd{is\_term\_candidate}: the first and last words cannot be function
    words (with/from/of\ldots), 3+ word phrases cannot have conjunctions/prepositions
    in the middle (and/or/with\ldots), and all-lowercase single words must be at least
    4 characters and contain an uppercase letter or hyphen.
  \item \textbf{TF-IDF calculation}---term frequency is the sum of candidate
    occurrences in each weighted segment multiplied by the segment weight; document
    frequency is the number of papers among the library's most recent 500 that contain
    the candidate. The IDF formula is $\ln\!\bigl((N+1)/(df+1)\bigr) + 1$. Three
    correction factors are applied:
    \begin{itemize}
      \item \textbf{section\_hits} ($\sqrt{\min(3, \text{number of sections appeared in})}$)
        ---appearing in more distinct sections increases credibility.
      \item \textbf{surface\_quality}---all-uppercase abbreviations $\times 2.0$,
        Title Case multi-word $\times 1.6$, hyphenated $\times 1.3$,
        all-lowercase multi-word $\times 0.7$. Surface forms that look more like
        technical terms score higher.
      \item \textbf{abbrev\_bonus} ($\times 1.8$)---candidates from the abbreviation-pair
        path receive extra weighting.
    \end{itemize}
  \item \textbf{Sort and truncate}---candidates with a total score $\geq 0.9$ are sorted
    in descending order and the top 12 are kept. Scores below the threshold are discarded
    (for example, a single-occurrence all-lowercase bigram scores about 0.7, just below
    the cutoff).
\end{enumerate}

\subsubsection{Abbreviation Pair Recognition}

For text like ``Structural Similarity Index Measure (SSIM),'' LitFolio uses a dedicated
regex to capture the full name and abbreviation, then performs two checks:

\begin{itemize}
  \item \textbf{Initial matching}---take the first letter of each word in the full name
    and greedily match against the abbreviation character by character; at least half
    must match. Hyphens are treated as word boundaries (``Retrieval-Augmented
    Generation'' $\to$ RAG matches).
  \item \textbf{Full-name refinement}---if the regex-captured full name is longer than
    the abbreviation, a word sequence from the tail equal to the abbreviation length
    is extracted for exact initial matching; if that misses, stop words (and/the/of)
    are removed and the process retries. For example, ``average structural similarity
    index measure (SSIM)'' is refined to ``structural similarity index measure.''
\end{itemize}

Only the abbreviation enters the final term list (the full name does not appear
independently), but the full name is saved as context for the LLM definition.

\subsubsection{Author Name Filtering}

A frequently overlooked noise source is author surnames---common surnames like ``Chen,''
``Yang,'' and ``Li'' appear at high frequency in paper author lists, satisfy Title Case
proper-noun criteria, and can rank highly among term candidates. Before extraction,
LitFolio collects all author surnames of the current paper into a HashSet (splitting
on spaces, hyphens, and periods, then lowercasing). Candidates that match after
normalization are skipped. This is especially important for Chinese authors' pinyin
surnames---``Chen'' and ``Wang'' appear at extremely high frequency in author lists
and would dominate the term rankings without filtering.

\subsubsection{LLM Definition Generation}

Once candidates are finalized, each term's ``surface form + full name (if any) +
evidence snippet'' is assembled into a prompt for the LLM, which returns a JSON-formatted
one-line Chinese definition. If the LLM call fails or returns empty, a fallback is used:
\cmd{This paper uses <term> in the context of the current argument or method.}

\begin{tipbox}
Term generation consumes significantly more tokens than TL;DR (it must scan the full
text). Consider binding it to a cost-effective domestic low-tier conversational model,
or running a local 7B+ open-source model.
\end{tipbox}

\section{Structured Note Cards}

The ``Notes'' tab has been upgraded from a single text box to a \textbf{structured card
system}. Each paper has five independently editable note sections by default, stored in
the \cmd{paper\_note\_sections} table.

\subsection{Default Sections}

\begin{description}
  \item[Problem] What problem the paper addresses.
  \item[Method] What method is proposed.
  \item[Key Numbers] Key experimental data or metrics.
  \item[Limitations] Limitations and open issues.
  \item[Thoughts] Your own ideas and comments.
\end{description}

Each section is independently editable, collapsible, and reorderable by drag-and-drop.
Sort order is persisted through the \cmd{sort\_order} field. You can add custom sections
(e.g., ``Experiment Ideas,'' ``Related Work''); custom sections have the same editing
and sorting capabilities as the defaults.

\subsection{AI Auto-Fill Mechanism}

Each section has a \cmd{source} field indicating the content origin:

\begin{description}
  \item[user] Content you wrote manually.
  \item[ai:tldr] Auto-filled by TL;DR (Quick Read).
  \item[ai:quick\_read] Auto-filled by Deep Read.
\end{description}

The fill rule is \textbf{fill only empty, never overwrite}:

\begin{enumerate}
  \item When TL;DR runs, the \cmd{key\_findings} array returned by the LLM is written
    to the ``Key Numbers'' section---but only if that section's \cmd{content} is empty
    and its \cmd{source} is not \cmd{user}.
  \item When Deep Read runs, the four-section JSON is mapped to the corresponding
    sections: \cmd{problem} $\to$ Problem, \cmd{method} $\to$ Method,
    \cmd{limitations} $\to$ Limitations. Only empty sections are filled.
  \item If you have manually edited a section (\cmd{source = user}), AI will never
    overwrite it, even if you re-run TL;DR or Deep Read.
  \item AI-filled content is prefixed with a timestamp: \cmd{[2025-05-27 AI:TLDR]
    <content>}.
\end{enumerate}

\subsection{Sections and Full-Text Search}

All sections' \cmd{content} fields are included in the \cmd{papers\_fts} FTS5 index.
This means you can search for a keyword you wrote in a section and find the
corresponding paper---even if the word never appears in the title or abstract.

\begin{tipbox}
Section drag-and-drop sorting uses HTML5 Drag and Drop; sort order is written to
the \cmd{sort\_order} field. The default order is Problem $\to$ Method $\to$ Key Numbers
$\to$ Limitations $\to$ Thoughts.
\end{tipbox}

\section{Annotation Color Semantics}

Highlight colors can carry semantic meaning, stored in the \cmd{highlights.label} field.
Configure color-to-label mappings under ``Settings $\to$ Annotation Colors.'' Default
mapping:

\begin{longtable}{p{2cm} p{3.5cm} p{7cm}}
\toprule
\textbf{Color} & \textbf{Label} & \textbf{Semantics} \\
\midrule
\endhead
Yellow & Key Finding & Key findings, important conclusions \\
Purple & Method & Methodological points, algorithmic details \\
Blue & To Verify & Claims that need verification, questionable data \\
Red & Disagree & Points of disagreement, controversial assumptions \\
Green & Citable & Excellent formulations worth quoting directly \\
\bottomrule
\end{longtable}

\subsection{Workflow}

\begin{enumerate}
  \item \textbf{Create highlights}---after selecting text, the popup menu shows each
    color with its label name (e.g., ``Yellow $\cdot$ Key Finding''). Choosing a color
    simultaneously sets the label.
  \item \textbf{Filter highlights}---the top of the reader's left pane now has a row
    of label filter chips. Clicking a chip shows only highlights with that label;
    click again to remove the filter. Multiple chips can be active (OR semantics).
  \item \textbf{AI integration}---the highlight summary feature (invoke LLM on selected
    highlights) supports label-based filtering. For example, summarize only ``Key
    Finding'' highlights, ignoring ``To Verify'' noise.
  \item \textbf{Export integration}---during Markdown export, highlights are grouped
    by label, with each label as a third-level heading.
\end{enumerate}

\subsection{Custom Labels}

You can modify the color-to-label mapping in Settings: rename labels, add new colors,
or delete unused mappings. Existing highlights are unaffected---the label name is stored
in each highlight record, not referenced from the global configuration.

\section{Markdown Notes}

The bottom of the ``Notes'' tab still has a free-text Markdown editing area, auto-saved
to \fpath{papers/<id>/note.md}. LitFolio does not enforce any template---you can use
any Markdown syntax and embed arbitrary image links; the file is synced along with
the PDF via WebDAV.

\section{PDF Search}

\kbd{Ctrl+F} (macOS: \kbd{Cmd+F}) opens the search box:

\begin{itemize}
  \item \kbd{Enter} jumps to the next match; \kbd{Shift+Enter} jumps to the previous.
  \item Match count is displayed to the right of the search box.
  \item \kbd{Esc} closes the search box.
\end{itemize}

\section{Page Shortcuts}

Other commonly used keys in the reader:

\begin{itemize}
  \item \kbd{j} / \kbd{k}---next page / previous page.
  \item \kbd{1} / \kbd{2} / \kbd{3}---switch the right-pane tab.
  \item \kbd{[} / \kbd{]}---adjust left/right pane width (5\% step).
\end{itemize}

The full shortcut table is in \textbf{Chapter 17}.
